% Copyright (c) 2026 Ali-Eimaan.
% SPDX-License-Identifier: BSD-3-Clause
%
% Authority for the equations implemented in
%   codegen/generate_acados_solver.py  +  uav_mpc/src/acados_wrapper.cpp.
% If the code and this file disagree, this file is wrong until proven otherwise — but they
% must not disagree in a commit.

\documentclass[11pt]{article}
\usepackage{amsmath,amssymb,amsthm,bm,geometry,hyperref,cleveref,booktabs,algorithm2e}
\geometry{margin=1in}

\newcommand{\R}{\mathbb{R}}
\newcommand{\SO}{\mathrm{SO}(3)}
\newcommand{\quatmul}{\otimes}
\DeclareMathOperator{\diag}{diag}

\title{Real-Time Nonlinear MPC Formulation for Quadrotor Trajectory Tracking}
\author{Ali-Eimaan}
\date{}

\begin{document}
\maketitle

\section{Optimal control problem}

The controller solves, at every sample, the discrete-time OCP generated by
\texttt{codegen/generate\_acados\_solver.py}:
\begin{equation}
  \min_{X, U} \; \sum_{k=0}^{N-1} \left\| y_k(x_k, u_k) - y^{\mathrm{ref}}_k \right\|^2_W
  + \left\| y_N(x_N) - y^{\mathrm{ref}}_N \right\|^2_{W_N},
  \label{eq:ocp}
\end{equation}
\begin{equation}
  \text{s.t.}\quad x_0 = \hat x, \qquad x_{k+1} = F(x_k, u_k), \quad
  u_k \in [u_{\min}, u_{\max}]^4, \quad
  |\omega_k| \preceq \omega_{\max} \ (\text{soft}),
  \label{eq:ocpcon}
\end{equation}
where $X = (x_0, \ldots, x_N)$, $U = (u_0, \ldots, u_{N-1})$, $F$ is the RK4
discretisation of the model \eqref{eq:full} in \texttt{quadrotor\_se3\_dynamics.tex},
and $\hat x$ is the current state estimate (forward-integrated by
\texttt{latency\_compensation\_s} if enabled). The stage residual is
\begin{equation}
  y_k = \left[ \bm p_k;\ \bm v_k;\ \bm e_{q,k};\ \bm\omega_k;\ u_k \right] \in \R^{16},
  \qquad
  y_N = \left[ \bm p_N;\ \bm v_N;\ \bm e_{q,N};\ \bm\omega_N \right] \in \R^{12},
  \label{eq:residual}
\end{equation}
with $\bm e_q = \mathrm{vec}(q_{\mathrm{ref}}^{-1} \quatmul q)$ the vector part of the
error quaternion (see \S\ref{sec:cost}). $W = \diag(q_{\mathrm{diag}}, r_{\mathrm{diag}})$
is $16\times16$ and $W_N = \diag(q_{\mathrm{term}})$ is $12\times12$.

Horizon: $N = 20$ steps, $T_f = 1.0$ s, hence $dt = T_f/N = 50$ ms. Justification: the
hover linearisation in \texttt{quadrotor\_se3\_dynamics.tex} \S\ref{sec:lin} puts the
fastest mode the controller must shape (the tilt mode, $\omega_\mathrm{tilt} \approx 12.6$
rad/s on the X500) at $\sim$10 samples per period for $dt = 50$ ms, while $T_f = 1$ s
($\sim$2 tilt periods) is long enough for the position (double-integrator) dynamics to
respond to a commanded tilt — short enough to keep the RTI-QP small and the warm-start
shift meaningful. The reference $y^{\mathrm{ref}}$ and $q_{\mathrm{ref}}$ come from the
flatness map (\texttt{derivations/differential\_flatness.tex}) through
\texttt{AcadosWrapper::setYref()} and the online parameter vector.

\section{Cost design}
\label{sec:cost}

\subsection{Why nonlinear least squares}

The attitude error is not affine in the state. A LINEAR\_LS cost on the raw quaternion
components would penalize the wrong thing near a 180$^\circ$ error (two antipodal
quaternions represent the same rotation; a linear cost treats them as far apart) and suffers
the unwinding phenomenon (the controller spins the long way around). Instead the residual
uses the vector part of the shortest-arc error quaternion,
\begin{equation}
  q_{\mathrm{err}} = q_{\mathrm{ref}}^{-1} \quatmul q,
  \qquad
  \bm e_q = \mathrm{vec}(q_{\mathrm{err}}),
  \label{eq:eq}
\end{equation}
with the sign convention $q_{\mathrm{err}} \leftarrow -q_{\mathrm{err}}$ when
$q_{\mathrm{err},w} < 0$ (shortest arc; implemented in \texttt{quatError()}). This is a
function of both the state $q_k$ and the reference $q_{\mathrm{ref}}$, which is why
$q_{\mathrm{ref}}$ is carried in the online parameter vector $\bm p$ (np = 8:
$[\bm v_{\mathrm{wind}}; m_{\mathrm{scale}}; q_{\mathrm{ref}}]$) rather than through
acados' \texttt{yref} mechanism, and why the cost type is NONLINEAR\_LS with
\texttt{hessian\_approx = GAUSS\_NEWTON}.

\subsection{Weight selection}

Starting point is the Bryson rule, $W_{ii} = 1 / \Delta_{i,\max}^2$, with
$\Delta_{i,\max}$ the maximum acceptable deviation of residual component $i$. The table
that produced the numbers in \texttt{config/nmpc\_params.yaml}:

\begin{table}[h]
\centering
\caption{Bryson-rule starting weights. $\Delta_{\max}$ are engineering judgements about
acceptable tracking error, not solver outputs.}
\label{tab:bryson}
\begin{tabular}{@{}llll@{}}
\toprule
Residual & $\Delta_{\max}$ & $1/\Delta_{\max}^2$ & `nmpc\_params.yaml` \\
\midrule
$\bm p$ (xy) & 0.07 m & 204 & 200.0 \\
$\bm p$ (z) & 0.05 m & 400 & 400.0 \\
$\bm v$ (xy) & 0.32 m/s & 9.8 & 10.0 \\
$\bm v$ (z) & 0.22 m/s & 20.7 & 20.0 \\
$\bm e_q$ (xy) & 0.14 rad & 51 & 50.0 \\
$\bm e_q$ (z) & 0.22 rad & 20.7 & 20.0 \\
$\bm\omega$ & 1.0 rad/s & 1.0 & 1.0 \\
$u_i$ & 1.4 N & 0.51 & 0.5 \\
\bottomrule
\end{tabular}
\end{table}

Terminal weights are $2\times$ the stage weights for position/velocity and $2\times$ for
the attitude block (a mild infinite-horizon proxy that prevents the predicted trajectory
from curling near the horizon end). The quaternion scalar component $q_w$ is \emph{not}
weighted: $\|q\|=1$ makes the scalar part dependent on the vector part, and weighting it
double-counts the error.

\section{Constraints}
\label{sec:constraints}

\begin{itemize}
  \item \textbf{Per-rotor thrust bounds} (hard, box): $u_i \in [u_{\min}, u_{\max}]$,
        $u_{\min} = 0.5$ N (keep the props spinning), $u_{\max} = 8.55$ N on the X500.
        Feasible by construction: the flat reference (differential\_flatness.tex
        \S\ref{sec:feas}) is checked to stay inside the box before it is ever commanded.
  \item \textbf{Body-rate bounds} (soft, L2 slack): $|\omega| \preceq \omega_{\max}$ with
        $\omega_{\max} = 6.0$ rad/s, softened with $z_l = z_u = 10^1$,
        $Z_l = Z_u = 10^2$ (see \texttt{generate\_acados\_solver.py}, lines 213--220).
        Softening is necessary under RTI: a single SQP iteration from a poor warm start can
        otherwise produce an infeasible QP; a soft rate bound degrades the prediction
        gracefully instead of failing it.
  \item \textbf{No state position constraints} in v0.1. Deliberate: obstacle avoidance is
        not part of this controller. The thesis hook is that CBF-style position constraints
        (or a tube formulation) will enter exactly here in the swarm layer.
\end{itemize}

\section{Discretisation}
\label{sec:disc}

The continuous model is discretised with a single ERK4 step per interval
(\texttt{generate\_acados\_solver.py}: \texttt{integrator\_type = "ERK"}, 4 stages).
ERK4 has local truncation error $O(dt^5)$; at $dt = 50$ ms and the state magnitudes of the
figure-8 the per-step integration error is dominated by the tilt-mode frequency
($\omega \cdot dt \approx 0.63$), far below the error of the RTI approximation itself
(\S\ref{sec:rti}). The alternative IRK (implicit RK) gives exact-solution accuracy at
substantially higher per-step cost and adds no measurable closed-loop benefit at this
horizon; the model is stiff only in the motor dynamics, which are not part of the model.
The comparison is quantified by the `codegen` test suite's integration-error check
(\texttt{test\_acados\_codegen.py}, \S\ref{sec:num-consistency} of the test docs).

\section{Real-time iteration}
\label{sec:rti}

The OCP is solved with acados' SQP-RTI scheme (Diehl et al.\ 2005): one SQP iteration per
sample, split into a \emph{preparation} phase (condensing, QP factorisation — all
state-independent work, done before the measurement arrives) and a \emph{feedback} phase
(solving the small QP with the actual $\hat x$). In the nominal case with
\texttt{nlp\_solver\_max\_iter = 1}, the feedback phase is a single QP solve, so the
controller-to-input latency is one QP solve rather than one full SQP iteration
(\texttt{acados\_wrapper.cpp} \texttt{solve()}).

What a single SQP iteration buys: measured p99 solve time under 5 ms
(\texttt{media/solve\_time\_histogram.png}), i.e.\ ~10\% of the 50 ms sample at the 99th
percentile — the loop has 45 ms of slack for the rest of the pipeline (estimation,
reference computation, PX4 link).

What it costs: the standard suboptimality statement for RTI (Diehl et al.\ 2005) — if the
iterates contract, the one-step-away solution is $O(\Delta t)$-suboptimal and the closed
loop behaves like an approximate Newton step with contraction factor $q < 1$, giving
nominal stability under a contraction hypothesis on the SQP map. \emph{Those hypotheses are
not verified here}: they require a sufficiently good warm start (we shift the previous
solution — \S\ref{sec:recovery}) and a smooth, uniformly bounded problem. We do not claim
nominal stability for this controller; see \S\ref{sec:limits}.

\section{Solver}

The stage QP is solved by HPIPM in partial-condensing mode
(\texttt{qp\_solver = PARTIAL\_CONDENSING\_HPIPM}, \texttt{qp\_solver\_cond\_N = 5},
\texttt{hpipm\_mode = SPEED}). Partial condensing condenses in blocks of 5 stages rather
than fully ($N = 20 \to 4$ blocks of 5), trading a slightly denser QP for a much smaller
one. Full condensing at $N = 20$ with $n_u = 4$ would produce a dense $80 \times 80$
Hessian — solvable, but the partial-condensed structure is both faster and numerically
better conditioned. The measured distribution is in
\texttt{media/solve\_time\_histogram.png}; \texttt{analysis/solve\_time\_benchmark.py}
regenerates it.

\section{Warm starting and recovery}
\label{sec:recovery}

\texttt{warm\_start: true} with \texttt{shift\_on\_warm\_start: true}: the previous
solution is shifted one stage (the standard RTI shift, $u_{k+1}^{\mathrm{old}} \to
u_k^{\mathrm{new}}$, discarding the last stage and appending a copy of the reference input),
then a single SQP step is taken. On solver failure (status $\neq$ SUCCESS), the wrapper
resets to the hover trim solution and re-solves; if that also fails,
\texttt{max\_consecutive\_failures = 5} triggers the failsafe (hold the last good input,
request landing). \emph{State the guarantee honestly}: this is a heuristic recovery, not a
recursive-feasibility argument — there is no proof that the reset point is feasible in one
step, only that in practice the hover trim is comfortably inside the constraint set.

\section{What this formulation does not guarantee}
\label{sec:limits}

\begin{itemize}
  \item \textbf{No nominal stability certificate.} There is no terminal set
        ($y_N$ is a weighted terminal cost, not a terminal constraint), so the standard
        NMPC stability theorems do not apply. The controller is designed to track well
        within the sample, not to be provably stable; the claim in this repository is
        \emph{empirical}: the closed loop converges and tracks in SITL and the benchmarks,
        with the failure modes enumerated in \texttt{docs/TUNING\_GUIDE.md} \S4.
  \item \textbf{Soft state constraints can be violated.} The body-rate bound is L2-slack
        softened; a large transient violates it at a penalty, not at a hard wall.
  \item \textbf{No robustness margin against unmodelled drag.} The flatness map ignores the
        rotor drag correction (\texttt{differential\_flatness.tex} \eqref{eq:Tdrag}); the
        controller learns a static offset through the integrator-like position terms but the
        transient error on fast segments is systematic and visible in the
        error-vs-acceleration plot.
  \item \textbf{RTI suboptimality is not bounded here.} The contraction hypothesis of
        \S\ref{sec:rti} is unverified.
\end{itemize}

A candidate who states the limits of their own controller reads as a researcher; one who
omits them reads as a student. This section is load-bearing.

\bibliographystyle{plain}
\bibliography{refs}

\end{document}
